\chapter{\textcolor{mycolor}{\textbf{État de l'art et fondements théoriques}}}
\markboth{\chaptername\ \thechapter.\ État de l'art et fondements théoriques}{}
\fancyhf{}
\fancyhead[L]{\small\itshape\nouppercase{\rightmark}}
\fancyfoot[L]{\small\itshape\nouppercase{\leftmark}}
\fancyfoot[R]{\small\itshape\thepage}
\renewcommand{\headrulewidth}{1pt}
\renewcommand{\footrulewidth}{0.4pt}

\vfill
\begin{flushright}{\Large\bfseries\textcolor{mycolor}{Introduction}}\end{flushright}
\addcontentsline{toc}{section}{Introduction}

Avant de présenter nos choix, nous situons le projet dans son domaine. Ce chapitre passe en revue les assistants personnels et l'informatique contextuelle, puis les quatre familles de technologies sur lesquelles repose la chaîne : l'Internet des objets et l'Edge AI, les données et leur évaluation, la vision par ordinateur et les modèles vision-langage, et les petits modèles de langage. Il se termine par une synthèse qui justifie l'approche retenue.

\clearpage

\section{Assistants personnels et informatique contextuelle}

\aRediger{6 (coordination)}{Présenter les assistants vocaux et applications de rappel existants, leurs limites (commandes explicites, absence de contexte) et les travaux sur l'assistance contextuelle.}

% Sous-équipe 1 : IoT et matériel. Section du chapitre 2, appelée par rapport/main chapters.
% Retirer l'encadré \aRediger une fois la section écrite.

\section{Internet des objets, cartes embarquées et Edge AI}

\aRediger{1 (IoT et matériel)}{Présenter les architectures IoT, les caméras et capteurs utilisables, le Raspberry Pi et les contraintes de l'Edge AI (mémoire, calcul, consommation).}

\subsection{Architectures IoT}

\subsection{Caméras et capteurs pour la détection de présence}

\subsection{Raspberry Pi et traitement embarqué}

\subsection{Edge AI et traitement centralisé}


% Sous-équipe 2 : Data. Section du chapitre 2, appelée par rapport/main chapters.
% Retirer l'encadré \aRediger une fois la section écrite.

\section{Données, annotation et évaluation}

\aRediger{2 (Data)}{Présenter les jeux de données existants, la méthode d'annotation retenue et les métriques qui serviront au chapitre 5.}

\subsection{Jeux de données de détection d'objets}

\subsection{Constitution et annotation d'un jeu de scénarios}

\subsection{Précision, rappel, F1-score et matrice de confusion}

\subsection{Mesure de la latence}


% Sous-équipe 3 : Vision (VLM). Section du chapitre 2, appelée par rapport/main chapters.
% Retirer l'encadré \aRediger une fois la section écrite.

\section{Vision par ordinateur et modèles vision-langage}

\aRediger{3 (Vision (VLM))}{Présenter les fondements de la détection d'objets, les détecteurs légers et les modèles vision-langage, puis comparer les options.}

\subsection{Détection et classification d'objets}

\subsection{La famille YOLO}

\subsection{Les modèles vision-langage : LLaVA et Moondream2}

\subsection{Comparaison des approches candidates}


% Sous-équipe 4 : Décision et backend. Section du chapitre 2, appelée par rapport/main chapters.
% Retirer l'encadré \aRediger une fois la section écrite.

\section{Petits modèles de langage et fusion multimodale}

\aRediger{4 (Décision et backend)}{Présenter les petits modèles de langage candidats, leur exécution locale (quantification, Ollama) et les approches de fusion.}

\subsection{Des grands aux petits modèles de langage}

\subsection{Phi, Qwen et Mistral}

\subsection{Quantification et exécution locale}

\subsection{Fusion de données multimodales et règles de décision}


\section{Synthèse et positionnement de LIFEWISE}

\aRediger{6 (coordination)}{Comparer les approches étudiées et justifier les choix de LIFEWISE pour le scénario pilote.}


\section*{Conclusion}
\addcontentsline{toc}{section}{Conclusion}

Ce chapitre a présenté les technologies disponibles pour chaque maillon de la chaîne. Le chapitre suivant décrit comment nous les assemblons dans l'architecture de LIFEWISE.
